\documentclass[11pt,a4paper]{article}
\usepackage[T2A]{fontenc}
\usepackage[utf8]{inputenc}
\usepackage[russian]{babel}
\usepackage{amsmath,amssymb,amsthm,mathtools}
\usepackage{geometry}
\geometry{margin=2.4cm}
\usepackage{booktabs}
\usepackage{enumitem}
\usepackage{caption}
\usepackage{titlesec}
\usepackage[colorlinks=true,linkcolor=black,urlcolor=blue,citecolor=black]{hyperref}

\theoremstyle{definition}
\newtheorem{definition}{Определение}[section]
\newtheorem{hypothesis}[definition]{Гипотеза}
\newtheorem{proposition}[definition]{Утверждение}
\newtheorem{criterion}[definition]{Критерий}
\newtheorem{remark}[definition]{Замечание}
\theoremstyle{plain}
\newtheorem{theorem}[definition]{Теорема}
\newtheorem{lemma}[definition]{Лемма}
\newtheorem{corollary}[definition]{Следствие}

\newcommand{\R}{\mathbb{R}}
\newcommand{\PP}{\mathcal{P}}
\newcommand{\sph}{\mathbb{S}}
\newcommand{\norm}[1]{\left\lVert #1 \right\rVert}
\newcommand{\abs}[1]{\left\lvert #1 \right\rvert}
\newcommand{\inner}[2]{\left\langle #1, #2 \right\rangle}
\DeclareMathOperator{\tr}{tr}
\DeclareMathOperator{\spec}{spec}
\DeclareMathOperator{\softmax}{softmax}
\DeclareMathOperator{\rank}{rank}

\title{\textbf{Представление понятия в трансформере:\\ программа исследования}\\[0.4em]
\large Постановка задачи, гипотеза непрерывности, конструкции и эксперименты}
\author{}
\date{}

\begin{document}
\maketitle

\begin{abstract}
\noindent
Документ фиксирует программу исследования, выросшую из полемики с теорией концептуальных
пространств П.~Гарденфорса. Исходная претензия: геометрический словарь в этой теории не несёт
за собой соответствующего математического аппарата — «выпуклость», «метрика», «домен»
употребляются без теорем и следствий.

Трансформер отличается тем, что его представления не постулируются, а возникают как следствие
оптимизации простой задачи: предсказать следующий токен. Отсюда единственный допустимый ход —
\emph{выводить} свойства представлений из того, что оптимизируется при обучении, а не назначать
геометрию заранее. Такой вывод даёт гипотезу, а не основание: механизм внимания обучается
кодировать связи, полезные для предсказания, и правдоподобно, что для текстов одного понятия
эти связи одни и те же. Это и есть предмет проверки.

Документ намеренно устроен как программа, а не как утверждение. Центральная трудность —
отсутствие эталона: чтобы проверить, близки ли представления текстов одного понятия, нужно
заранее знать, что это «одно понятие», а это и есть искомое. Разбирается, как разорвать этот
круг измеримым способом: какие инварианты у слоя вообще существуют (не значения векторов и не
их число, а соотношения между ними), что из них следует из обучения, и какие эксперименты
отличают гипотезу от артефакта. Отдельно разбирается, что каждый эксперимент \emph{не}
позволяет заключить: главный порок обсуждаемой литературы — подмена проверяемых утверждений
метафорами.
\end{abstract}

\tableofcontents
\bigskip

%==========================================================
\section{Постановка проблемы}

\subsection{Исходное возражение}

Теория концептуальных пространств моделирует понятие как выпуклое подмножество пространства
перцептивных качеств, а сходство — как расстояние в этом пространстве. Возражение состоит в
том, что термины здесь опережают аппарат. За словом «тензор» стоит тензорное исчисление с
законами преобразования и свёртками; за «выпуклостью» — выпуклый анализ с теоремами
отделимости и двойственности. Теория концептуальных пространств не использует ни того, ни
другого: выпуклость в ней — проверяемое эмпирическое допущение, а не следствие аксиоматики,
метрика постулируется, а не выводится из структуры качества, а классическая критика
(А.~Тверски) показывает, что человеческое сходство вообще не удовлетворяет аксиомам метрики —
симметрии и неравенства треугольника. Тем самым геометрический образ является
\emph{выбранным языком описания}, а не установленной структурой.

Отсюда не следует, что геометрия неприменима в принципе. Следует более узкое и более
сильное требование: геометрическая модель допустима там, где порождающая её структура
задана независимо и допускает проверку. Ниже разбирается, даёт ли трансформер такую структуру.

\subsection{Что известно по факту, а что назначается}

В трансформере есть текст, есть токены и есть векторы на выходе слоёв. Это — \emph{факт}:
матрица $V^{(\ell)}\in\R^{n\times d}$ на каждом слое $\ell$, где $n$ — число позиций, $d$ —
размерность канала. Никакое понятие в этой матрице не лежит «готовым»: понятие — это то, что
мы приписываем тексту, наблюдая его векторы.

Из этого следует важное ограничение. Нельзя объявить, что представление понятия — это, скажем,
пул векторов или их главное направление, и затем доказывать устойчивость этого пула. Такой
ход защищает \emph{назначенную} конструкцию, но не отвечает на вопрос, отражает ли она хоть
что-то. Именно этим недостатком страдает исходная теория концептуальных пространств, и
повторить его, поменяв язык с «перцептивных качеств» на «векторы слоя», было бы бессмысленно.

Единственный допустимый ход — идти от того, что \emph{оптимизируется при обучении}. Не
назначать геометрию, а вывести, какие свойства представлений обязаны возникнуть из задачи
оптимизации, и затем проверить, возникли ли они.

\subsection{Что оптимизирует attention и что из этого следует}

Обучение идёт на предсказании следующего токена, функция потерь — кросс-энтропия
\[
\mathcal{L}=-\sum_t \ln P(x_{t+1}\mid x_1,\dots,x_t).
\]
Механизм внимания на каждом слое вычисляет
\begin{equation}\label{eq:attn}
A_{ij}=\softmax_j\!\left(\frac{q_i^\top k_j}{\sqrt d}\right),\qquad
\mathrm{out}_i=\sum_j A_{ij}v_j ,
\end{equation}
где $q_i=W_Q x_i$, $k_j=W_K x_j$, $v_j=W_V x_j$ — проекции входного вектора. Обучение
подстраивает $W_Q,W_K,W_V$ так, чтобы уменьшить $\mathcal{L}$.

Из этого выводится следующее, и здесь необходимо разделять, что следует строго, а что
является правдоподобной гипотезой.

\begin{proposition}[Что следует строго]\label{prop:strict}
Матрица $A$ не является мерой семантической близости. Она является распределением, обученным
так, чтобы взвешенная сумма значений $\sum_j A_{ij}v_j$ была полезна для предсказания
следующего токена. Скалярное произведение $q_i^\top k_j$ обучается как мера
\emph{полезности связи} для предсказания, а не как мера сходства содержаний.
\end{proposition}

Это утверждение — по существу, определение того, что оптимизируется, — а не эмпирическая
гипотеза. Из него, однако, \emph{не} следует автоматически, что тексты одного понятия дают
похожие представления. Вывести это можно только предположительно.

\begin{hypothesis}[Единство связей внутри понятия]\label{hyp:training}
Если два текста задают одно и то же понятие, то связи, полезные для предсказания внутри этих
текстов, в существенной части совпадают. Поэтому структура внимания $A$ у таких текстов
оказывается близкой — даже если сами векторы $V$ разнонаправлены и имеют разную длину.
\end{hypothesis}

Почему это правдоподобно: предсказание в описании пчелы опирается на одни и те же отношения —
«летает», «собирает», «живёт в сотах», «подчиняется королеве», — какую бы форму ни принимало
описание. Если модель выучила эти отношения как полезные, матрица внимания отразит их форму
независимо от конкретных слов.

Почему это не факт: функция потерь не содержит требования «тексты одного понятия дают близкие
$A$». Она требует только предсказуемости. Близость может возникнуть как побочный продукт, а
может и не возникнуть: модель вправе опираться на разные признаки для предсказания в разных
описаниях, и тогда гипотеза ложна. Проверить это можно только экспериментом
(раздел~\ref{sec:experiments}).

\subsection{Наблюдаемая трудность: контрпример}

Простая иллюстрация того, почему нельзя идти от направлений векторов. Рассмотрим три текста,
задающие одно понятие:
\begin{quote}
«пчела»; \quad «насекомое, собирающее мёд»; \quad «не птица, но летает, живёт в сотах и
подчиняется королеве».
\end{quote}
Векторы этих текстов на выходе слоёв направлены существенно по-разному: словарь и синтаксис
различны, и косинусы между векторами позиций малы. Если представление понятия строить как
направление (пул, взвешенный вниманием, или главный собственный вектор), эти три текста
разойдутся, хотя задают одно понятие. Отображение, назначившее представлением направление,
будет опровергнуто этим примером.

Этот контрпример и определяет, что искать: инвариант должен опираться не на направления
отдельных векторов, а на \emph{соотношения} между ними, сохраняющиеся от текста к тексту
(раздел~\ref{sec:invariants}).

\subsection{Предмет исследования}

Рассматриваются связные тексты $T$, задающие понятие: последовательность токенов
$x_1,\dots,x_n$, которой соответствуют векторы в промежуточных слоях и на выходе. Из текста
возникает не одна точка, а множество векторов $V=\{v_1,\dots,v_n\}$, из которых не все близки,
но которые совместно нечто задают. Требуется:
\begin{enumerate}[label=(\arabic*)]
\item определить, что из данных слоя (матрица $V$ и матрица $A$) является кандидатом в
инвариант понятия, обосновав это связью с задачей обучения, а не удобством;
\item определить, как сравнить эти инварианты у разных текстов \emph{без} заранее известного
эталона «одно понятие / разные понятия» (раздел~\ref{sec:nogroundtruth});
\item указать эксперименты, которые отличают подлинный инвариант от артефакта слоя, и явно
перечислить, что каждый из них не проверяет.
\end{enumerate}

\subsection{Архитектурный механизм: почему представление реляционно}\label{sec:mechanism}

Контрпример с пчелой (раздел~1.4) показывает, что искать надо в соотношениях, но не объясняет,
почему их следует ожидать, а не считать удачей. Объяснение даёт сама архитектура.

\begin{proposition}[Смешивание]\label{prop:mixing}
Пусть логиты $q_i^\top k_j/\sqrt d$ конечны. Тогда $\softmax$ строго положителен,
$A_{ij}>0$, и выход~\eqref{eq:attn} есть \emph{строгая} выпуклая комбинация всех векторов
значений:
\[
\mathrm{out}_i=\sum_{j=1}^{n}A_{ij}v_j,\qquad A_{ij}>0,\quad \sum_j A_{ij}=1 .
\]
В частности, $\mathrm{out}_i$ лежит во внутренности выпуклой оболочки $\{v_1,\dots,v_n\}$ и не
является функцией одного $x_i$.
\end{proposition}
\begin{proof}
$\softmax_j(z)=e^{z_j}/\sum_k e^{z_k}>0$ при конечных $z$. Сумма весов равна единице по
определению $\softmax$; из положительности всех слагаемых следует, что точка лежит во
внутренности оболочки. Зависимость от $x_i$ опосредована всеми $x_j$, поскольку $A_{ij}>0$.
\end{proof}

\begin{proposition}[Единственный межпозиционный канал]\label{prop:channel}
В слое трансформера MLP и LayerNorm действуют попозиционно. Поэтому зависимость
$v_i^{(\ell+1)}$ от $x_j^{(\ell)}$ при $j\ne i$ реализуется исключительно через матрицу
внимания $A^{(\ell)}$.
\end{proposition}

Из~\ref{prop:mixing} и~\ref{prop:channel} следует, что «изолированного» вектора-носителя
понятия в выходе слоя нет: контекст уже вложен в каждую позицию, а канал этого вложения —
отношение $A$. Содержательная информация о том, \emph{как} интерпретировать токен, имеет
реляционную природу (пара «кто на кого смотрит и с каким весом»), а не вид отдельной точки.

Почему тогда именно соотношения, а не направления? Остаточный поток между подслоями не имеет
выделенного базиса: линейные преобразования переопределяют координаты, и инвариантный смысл
несут лишь величины, построенные из скалярных произведений, — $\langle q_i,k_j\rangle$,
$\Sigma=V^\top V$, их спектры. Индивидуальные направления (координаты) не каноничны и как
носители смысла не определены. Эта калибровочная свобода \emph{приближённа}: попозиционные
нелинейности и LayerNorm ломают её, и насколько сильно направления всё же информативны —
эмпирический вопрос (его и ловит контрольная линия $\Phi_{\mathrm{pool}}$).

Сравнение со свёрточной сетью проясняет слово «нематериально». В свёртке правило смешивания
(ядро) фиксировано, разделяемо и потому визуализируемо: структура «вшита» в веса. В
самовнимании правило смешивания вычисляется из самого входа, $A=A(X)$, и не существует как
фиксированный печатаемый объект: оно реализуется лишь как отношение \emph{на данной}
конфигурации. «Нематериальность» означает здесь не отсутствие носителя, а три вещи: структура
распределена по позициям, зависит от экземпляра и определена с точностью до симметрий
(поворотов базиса). Именно поэтому её естественно искать в фактор-пространстве соотношений.

\begin{hypothesis}[Реляционное кодирование]\label{hyp:relational}
Понятийно-релевантная информация, выучиваемая при предсказании следующего токена (или
восстановлении зашумлённого входа), локализована в реляционных инвариантах — спектре
$\tilde\Sigma$ и структуре $A$, — а не в направлениях отдельных векторов; эти инварианты
совпадают у текстов одного понятия и устойчивы по слоям.
\end{hypothesis}

Это усиление гипотезы~\ref{hyp:training}: последняя утверждает совпадение связей, эта — что
связи и есть основной носитель смысла, а направления вторичны. Обе проверяемы
(раздел~\ref{sec:experiments}) и обе могут быть ложны; архитектурное рассуждение даёт лишь
необходимое условие правдоподобия, но не доказательство. Что именно проверяет каждый
эксперимент и чего он не проверяет — в разделе~\ref{sec:limits}.

%==========================================================
\section{Кандидаты в инвариант}\label{sec:invariants}

\subsection{Что меняется, а что может сохраняться}

Векторы слоя строятся динамически от слоя к слою и зависят от контекста. Поэтому значения
векторов, их направление и их число не подходят на роль инварианта понятия: они меняются при
переходе к другой формулировке, к другой длине текста и к другому слою. Кандидатом может быть
только то, что сохраняется от текста к тексту и от слоя к слою, — соотношения между векторами.

\subsection{Матрица Грама как минимальный инвариант соотношений}

\begin{definition}[Матрица Грама]
Пусть $V\in\R^{n\times d}$ — матрица векторов позиций. Матрицей Грама называется
\[
\Sigma=V^\top V\in\R^{d\times d},\qquad
\Sigma_{ab}=\sum_{i=1}^{n} v_{ia}v_{ib}=\sum_{i=1}^{n} (v_i)_a (v_i)_b .
\]
То есть $\Sigma=\sum_{i=1}^{n}v_i^\top v_i$ — сумма внешних произведений векторов позиций.
\end{definition}

Содержательно $\Sigma$ — это матрица всех попарных координатных соотношений векторов набора:
в ней нет ничего, кроме того, в каком отношении друг к другу находятся векторы. Её
собственные значения $\lambda_1\ge\dots\ge\lambda_d\ge 0$ показывают, сколько независимых
направлений реально занято набором.

Матрица Грама обладает ровно тем свойством, которое требуется по разделу~1.4: если весь набор
векторов повернуть одновременно ($V\mapsto RV$ для ортогональной $R$), то
\[
(RV)^\top(RV)=V^\top R^\top R\,V=V^\top V=\Sigma .
\]
Значит разнонаправленность векторов, о которой говорит контрпример с пчелой, $\Sigma$ не
разрушает. Это не назначение, а прямое следствие того, что $\Sigma$ хранит соотношения, а не
направления.

\begin{remark}[Что $\Sigma$ \emph{не} различает]
Два обстоятельства необходимо назвать сразу.
\begin{enumerate}[label=(\alph*)]
\item $\Sigma$ не различает повёрнутые друг относительно друга наборы векторов. Для гипотезы
это скорее необходимое свойство, но оно означает, что всякая информация, закодированная
только в направлении (а не в соотношениях), для $\Sigma$ потеряна.
\item $\Sigma$ не различает \emph{порядок} позиций: перестановка токенов не меняет сумму
внешних произведений. Между тем «собака кусает человека» и «человек кусает собаку» — разные
тексты. Значит $\Sigma$ в одиночку недостаточно; недостающая информация о порядке лежит в
матрице внимания $A$ (раздел~\ref{sec:order}).
\end{enumerate}
\end{remark}

\subsection{Нормировка и эффективный ранг}

Матрица $\Sigma$ растёт с числом позиций $n$ (каждое слагаемое добавляет вклад), поэтому
непосредственно сравнивать $\Sigma$ у текстов разной длины нельзя. Нормировка на след
$\tilde\Sigma=\Sigma/\tr\Sigma$ делает $\tilde\Sigma$ точкой симплекса плотностных матриц
$\{\rho\succeq0,\ \tr\rho=1\}$, где определены корректные информационные расстояния
(квантовая относительная энтропия и её симметризованная версия).

Спектр $\tilde\Sigma$ даёт числовой инвариант, не зависящий ни от длины, ни от порядка
нумерации координат:
\[
R(\Sigma)=\exp\!\Bigl(-\sum_k p_k\ln p_k\Bigr),\qquad p_k=\frac{\lambda_k}{\sum_j\lambda_j},
\]
--- эффективный ранг. $R$ показывает, в скольких направлениях реально «размазан» набор: $R=1$
означает, что все векторы сонаправлены (одно слово — вырожденный случай), большое $R$ —
набор растопырен по многим направлениям. Это кандидат в меру конкретности понятия, и его
поведение с ростом длины описания — проверяемое утверждение (раздел~\ref{sec:experiments}).

\subsection{Порядок: матрица внимания}\label{sec:order}

Матрица $A\in\R^{n\times n}$ хранит то, чего нет в $\Sigma$: какая позиция на какую смотрит.
Она несимметрична ($A_{ij}\ne A_{ji}$ в общем случае), и именно асимметрия кодирует порядок и
роли в предложении. Согласно гипотезе~\ref{hyp:training}, именно $A$ (или её спектр) должна
быть близка у текстов одного понятия, а $\Sigma$ — кодировать геометрию этих связей.

Трудность, которую необходимо назвать: $A$ имеет размер $n\times n$, а $n$ у сравниваемых
текстов разное («пчела» — один токен, развёрнутое описание — десятки). Сравнивать матрицы
разного размера напрямую нельзя. Возможные канонизации:
\begin{enumerate}[label=(\arabic*)]
\item \textbf{Спектр $A$.} Сравнивать наборы собственных значений (или сингулярных значений)
$A$. Число-инвариант, но, как и $\Sigma$, теряет порядок.
\item \textbf{Усреднение по позициям.} Свести $A$ к вектору весов позиций (например, диагонали
или строковым суммам). Грубо, но устойчиво к длине.
\item \textbf{Сопоставимая длина.} Сравнивать только тексты близкой длины. Честно, но резко
сужает корпус.
\end{enumerate}
Выбор канонизации — открытый вопрос программы, а не решение. Назвать его решённым было бы
нечестно: без него сравнение $A$ у текстов разной длины не определено.

%==========================================================
\section{Проблема отсутствия эталона}\label{sec:nogroundtruth}

Центральная методологическая трудность. Чтобы проверить гипотезу «тексты одного понятия дают
близкие инварианты», нужно знать, какие тексты задают одно понятие. Но именно это и есть
искомое: если бы у нас был критерий «одно понятие», гипотезу не нужно было бы проверять. Круг
замыкается.

Разорвать его можно, но это нужно делать явно и с оговорками.

\subsection{Ручной отбор как первичный эталон}

Простейший способ: взять тексты, про которые \emph{мы} уверены, что они задают одно понятие
(экспертная разметка), и проверить на них близость инвариантов. Это слабый, но законный
эталон при условии, что:
\begin{enumerate}[label=(\roman*)]
\item разметка делается \emph{до} вычисления инвариантов, чтобы не подгонять;
\item набор разнообразен (разные формулировки, разная длина), иначе гипотеза проверяется
только на удобных случаях;
\item отдельно держится контрольная группа «разные понятия», чтобы мерить не только
внутрипонятийную близость, но и межпонятийную отделимость.
\end{enumerate}
Главный изъян: экспертная разметка — это ещё одно мнение, и она не гарантирует, что
«одно понятие» понимается так же, как его понимает модель. Положительный результат на таком
наборе показывает лишь согласованность с нашей интуицией.

\subsection{LLM как разметчик, а не как судья}

LLM может расширить корпус: дать пары текстов с оценкой «одно понятие / разные». Это
разумно, но роль LLM необходимо ограничить.

\begin{remark}[Разделение ролей]
LLM допустима как \emph{разметчик}: она присваивает парам метки, по которым строится корпус.
После разметки LLM \emph{выбывает}: все дальнейшие вычисления — расстояния между инвариантами,
корреляции с метками, пороги отделимости — делаются формально, без её участия. Если же
спрашивать LLM «похожи ли эти тексты?» и сверять её ответ с нашей метрикой как эталон, то мы
сравниваем одну нейросеть с другой, обученной на тех же текстах. Это не проверка реальности, а
согласование двух сетей и может разделять общий дефект.
\end{remark}

Даже в роли разметчика LLM не даёт независимой истины: она обучена на тех же текстах, что и
исследуемая модель. Поэтому корпус, размеченный LLM, годится для \emph{расширения} и
\emph{проверки устойчивости} результата, но не как основание. Основание — ручной,
экспертный набор, а LLM-корпус — контроль, что результат не случаен на трёх примерах.

\subsection{Беспристрастный тест без эталона}

Есть тест, который не требует знания, какие тексты задают одно понятие. Если инвариант
действительно отражает понятие, он должен обладать \emph{устойчивостью по слоям}: у текстов
одного понятия эволюция инварианта от слоя к слою похожа, у текстов разных — расходится. Для
этого теста достаточно ручного набора, и он проверяет свойства, которые LLM не может знать
заранее.

Именно он ставится первым (раздел~\ref{sec:experiments}, шаг~1): он дешевле, честнее и
не требует ни эталона, ни LLM. Если он не даёт сигнала, никакая последующая калибровка корпуса
его не спасёт.

\subsection{Логика проверки: H1, H2 и положительный контроль}\label{sec:logic}

Прежде чем сравнивать инвариант с эталоном, полезно расщепить гипотезу на две, разные по
статусу проверяемости.

\begin{itemize}[label=(\roman*)]
\item \textbf{H1 (существование и надёжность).} Отображение $\Phi$ содержит воспроизводимый
сигнал: он не является шумом и не объясняется посторонними причинами (длина, частотность,
синтаксис, случайность весов).
\item \textbf{H2 (выравнивание с человеческими понятиями).} Индуцированная $\Phi$ структура
совпадает по соотношениям с нашей.
\end{itemize}

H1 проверяется \emph{без эталона}: это утверждение о приборе, а не о понятиях. H2 требует
независимого — человеческого — эталона. Смешение уровней приводит к неинтерпретируемым
результатам: отрицательный исход на уровне H2 неотличим от «$\Phi$ плохая» или «не тот
слой», если H1 не установлена.

\paragraph{Положительный контроль.} Чтобы отрицательный результат был содержательным, конвейер
сначала должен пройти валидацию: восстановить заведомо присутствующую структуру (синтетический
корпус с зашитыми соотношениями или известное реальное различение — число, согласование,
часть речи). Если $\Phi$ не видит даже этого, вывод об отсутствии понятийного сигнала
недействителен: сломан прибор, а не гипотеза.

\paragraph{Надёжность.} «Воспроизводимость» измеряется как корреляция двух независимых
измерений одной величины: расщепление выборки надвое (поправка Спирмена—Брауна), репликация по
независимо обученным моделям/архитектурам. Корреляцию с человеческим эталоном следует судить
относительно достижимого потолка (обе матрицы зашумлены), а не относительно нуля.

\paragraph{Выравнивание без круга.} H2 проверяется как \emph{изоморфизм структур второго
порядка}: строится матрица попарных расстояний между инвариантами понятий и независимая матрица
человеческих оценок сходства; сравниваются ранговой корреляцией. Это тест \emph{соотношений},
инвариантный к повороту и монотонному преобразованию, и он не требует совпадения самих точек.
Предсказательная валидность на отложенной части (held-out) отличает это от согласования с нашей
интуицией на том же корпусе.

Практические требования (контроль конфаундов, untrained-контроль, перемешивание порядка,
перестановочные тесты, поправка на перебор слоёв и метрик, защита от двойного использования
данных) сведены в шаге~0 раздела~\ref{sec:experiments}. Итог: строго проверяемо не тождество
понятий, а градуированное выравнивание структур с локализацией рассогласований.

%==========================================================
\section{Формальная постановка}

\subsection{Данные слоя}

\begin{definition}[Слой]
Для слоя $\ell=0,1,\dots,L$ обозначим через $V^{(\ell)}\in\R^{n\times d}$ матрицу векторов
позиций (по строке на позицию, по столбцу на канал), через $A^{(\ell)}\in\R^{n\times n}$ —
матрицу внимания, $A^{(\ell)}_{ij}\ge 0$, $\sum_j A^{(\ell)}_{ij}=1$, усреднённую по головам
или взятую по одной голове (выбор фиксируется до начала эксперимента).
\end{definition}

\subsection{Понятие как объект}

\begin{definition}[Отображение представления]
Отображением представления называется функция
\[
\Phi\colon \{\text{последовательности векторов}+\text{структура внимания}\}\longrightarrow \mathcal{M},
\]
где $\mathcal{M}$ — целевое пространство представлений, \emph{снабжённое метрикой}. Выбор
$\mathcal{M}$ не произволен: по разделу~\ref{sec:invariants} в него входят только функции от
соотношений ($\Sigma$, $A$), но не от направлений векторов. Для текста $T$ точка
$\Phi(T)\in\mathcal{M}$ называется \emph{представлением понятия, заданного текстом $T$};
сравнение двух текстов есть расстояние между их точками.
\end{definition}

\begin{definition}[Остров]
Пусть $\Phi$ фиксировано. Островом понятия называется множество текстов
\[
O_U=\{\,T:\ \Phi(T)\in U\,\},\qquad U\subset\mathcal{M}\ \text{область}.
\]
\end{definition}

\begin{remark}
Определение сознательно отделяет «остров» от «кластера в облаке токенов». Кластер токенов —
статистическая структура внутри одного текста. Остров — множество \emph{текстов}, чьи
представления близки. Смешение этих объектов делает постановку бессмысленной: компоненты
связности по попарным косинусам между токенами теряют порядок последовательности и
асимметрию внимания, то есть ровно ту информацию, которая отличает «собака кусает человека»
от «человек кусает собаку».
\end{remark}

\subsection{Требования к $\Phi$}

Отображение представления должно обладать тремя свойствами.

\begin{enumerate}[label=(\roman*)]
\item \textbf{Сжатие.} $\Phi$ переводит последовательность в один объект; без этого «остров»
негде расположить.
\item \textbf{Чувствительность к порядку и асимметрии.} $\Phi$ не должно быть инвариантно
относительно перестановки позиций и транспонирования ролей в предложении — иначе оно
описывает мешок слов, а не текст.
\item \textbf{Непрерывность.} Малое изменение входных векторов и матрицы внимания должно давать
малое изменение $\Phi$. Это техническое требование, и оно доказуемо (раздел~\ref{sec:stability}).
\end{enumerate}

%==========================================================
\section{Гипотеза непрерывности}

\begin{hypothesis}[Непрерывность представления]\label{hyp:continuity}
Пусть $T$ и $T'$ — два текста, задающих одно и то же понятие, и пусть $T'$ получен из $T$
малым изменением (перефразированием, добавлением уточняющего оборота). Тогда
\[
\norm{\Phi(T)-\Phi(T')}\ \le\ C\cdot \operatorname{dist}(T,T'),
\]
где $\operatorname{dist}$ — мера различия текстов, а $C$ — константа, не зависящая от текста.
Важно, что $\operatorname{dist}$ задаётся на уровне \emph{текста} (правка символов, замена
синонима), а не на уровне векторов: векторы — уже результат работы модели, и мерить ими
различие двух текстов означало бы снова опираться на то, что требуется найти. Какой именно
метрикой на текстах пользоваться — открытый вопрос; допустимый минимум — редакционное
расстояние плюс явно оговорённая близость по смыслу.
\end{hypothesis}

Содержательное усиление, которое и составляет суть программы:
\begin{hypothesis}[Отделимость понятий]\label{hyp:separation}
Для любых двух различных понятий $P\ne Q$ и любых текстов $T\in\mathcal{T}_P$, $S\in\mathcal{T}_Q$
величина $\norm{\Phi(T)-\Phi(S)}$ отделена от нуля снизу, причём типичное внутрипонятийное
расстояние существенно меньше типичного межпонятийного.
\end{hypothesis}

Вместе эти две гипотезы утверждают: тексты одного понятия стягиваются в компактную область,
отделённую от областей других понятий. Это и есть «остров» в точном смысле, и — что важнее
всего — это \emph{проверяемое} утверждение, а не философское.

\begin{remark}[Об области определения гипотезы]
Гипотеза формулируется для текстов, задающих понятие. Тексты, не задающие понятия (нарратив,
перечисление, риторическая фигура), лежат вне области её применимости, и их поведение
программой не покрывается. Это ограничение не техническое, а принципиальное: расширение на
все тексты потребовало бы отдельного критерия «задаёт понятие», которого пока нет.
\end{remark}

%==========================================================
\section{Конструкции отображения $\Phi$}

Предлагаются три конструкции. $\Phi_{\mathrm{pool}}$ служит базовой линией и — как показывает
контрпример из раздела~1.4 — заведомо опирается на неверный признак; $\Phi_{\mathrm{spec}}$
строится на соотношениях, инвариантных к направлению; $\Phi_{\mathrm{full}}$ добавляет
информацию о связях.

\subsection{Базовая линия: пул с нормировкой}

\begin{equation}
\Phi_{\mathrm{pool}}(T)=\frac{\sum_{i=1}^{n} w_i v_i}{\norm{\sum_{i=1}^{n} w_i v_i}},
\qquad w_i\ \text{--- вес позиции (равномерный или по маске).}
\end{equation}
Достоинство — прозрачность и дешевизна. Недостаток принципиален: отображение коммутативно по
позициям, и его значение целиком опирается на направление векторов. Контрпример с пчелой
показывает, что такая конструкция обязана расходиться на текстах одного понятия.
$\Phi_{\mathrm{pool}}$ сохраняется \emph{именно как контроль}: если она даёт компактные
области там, где конструкция на соотношениях — нет, значит направление в этом месте всё-таки
информативно; если расходится, как предсказано, — это подтверждение того, что искать нужно не
в направлении. Предварительный прогон (раздел~\ref{sec:prelim}) этого предсказания не
подтвердил: на малой модели направление, наоборот, несло основной сигнал.

\subsection{Конструкция на соотношениях: спектр матрицы Грама}

\begin{equation}
\Phi_{\mathrm{spec}}(T)=\spec(\tilde\Sigma),
\qquad \tilde\Sigma=\frac{V^\top V}{\tr(V^\top V)} .
\end{equation}
Здесь берётся \emph{только спектр} нормированной матрицы Грама, без всякого направления.
Конструкция инвариантна к повороту всего набора векторов — то есть к той разнонаправленности,
которую даёт контрпример, — и к длине текста (нормировка на след). Значения
$\Phi_{\mathrm{spec}}$ лежат в симплексе собственных значений, на котором определено расстояние
(например, квантовая относительная энтропия).

\paragraph{Вычислимая форма и метрика.} Собственные числа $V^\top V$ суть квадраты сингулярных
чисел $V$, поэтому
\[
\Phi_{\mathrm{spec}}(T)=p,\qquad p_i=\frac{\sigma_i^2(V)}{\sum_{j}\sigma_j^2(V)},
\]
где $\sigma_i(V)$ — сингулярные числа $V$ (убывающие). Это вероятностный вектор,
$p\in\Delta^{d-1}$, $\sum_i p_i=1$; при $n<d$ у него не более $\min(n,d)$ ненулевых компонент,
которые дополняются нулями до размерности $d$. Сравнение текстов разной длины корректно именно
потому, что все $p$ лежат в одном симплексе: расстояние берётся между распределениями —
например, Hellinger $d_H(p,q)=\sqrt{1-\sum_k\sqrt{p_kq_k}}$ или $W_1$ по отсортированным
компонентам. Один вызов сингулярного разложения даёт $\Phi_{\mathrm{spec}}$ без каких-либо
назначаемых параметров.

\paragraph{Что эта конструкция теряет.} Порядок позиций: как отмечено в
разделе~\ref{sec:invariants}, перестановка токенов не меняет $\Sigma$. Это ограничение
названо, а не скрыто; для учёта порядка вводится следующая конструкция.

\subsection{Учёт порядка: инвариант, использующий $A$}

Полный инвариант опирается на то, что $A$ кодирует связи, а $\Sigma$ — геометрию этих связей:
\begin{equation}
\Phi_{\mathrm{full}}(T)=\bigl(\spec(A),\ \spec(\tilde\Sigma)\bigr).
\end{equation}
Сингулярные значения $A$ отражают форму графа связей и инвариантны к перестановке нумерации
позиций. \emph{Однако направление связей они не различают}: у $A$ и $A^\top$ совпадают и
сингулярные, и собственные значения (характеристические многочлены совпадают), поэтому
$\spec(A)$ не отделяет «собака кусает человека» от «человек кусает собаку». Спектр хранит
лишь форму связей, но не их ориентацию. Для учёта направления нужен неспектральный дескриптор
$A$ (например, статистики асимметрии $A-A^\top$), и его выбор при разной длине текста остаётся
открытой проблемой (раздел~\ref{sec:order}).

\begin{remark}[Об отвергнутых конструкциях]
В более ранней редакции программы использовалась $\Phi_{\mathrm{attn}}$ — пул, взвешенный
вниманием, — и пара «направление, спектр». Контрпример из раздела~1.4 их опровергает: обе
опираются на направление векторов, которое у текстов одного понятия различно. Доказанные для
них оценки устойчивости остаются верными, но доказывают устойчивость \emph{не той}
конструкции. Теорема~\ref{thm:spec} об устойчивости спектра Грама сохраняет силу и относится к
$\Phi_{\mathrm{spec}}$; теорема~\ref{thm:attn} об устойчивости $A$ сохраняет силу и относится к
спектральной части $\Phi_{\mathrm{full}}$.
\end{remark}

\begin{remark}[Об эффективном ранге]
Величина
\[
R(\Sigma)=\exp\!\Bigl(-\sum_k p_k\ln p_k\Bigr),\qquad p_k=\frac{\lambda_k}{\sum_j\lambda_j},
\]
называется эффективным рангом и служит кандидатом в меру \emph{конкретности}: чем меньше $R$, тем
конкретнее понятие. Поскольку $\rank\Sigma\le \min(n,d)$, для коротких текстов $R$ ограничен
сверху числом позиций $n$; это ограничение учитывается при интерпретации. Сам по себе $R$ — не
представление, а одна из его координат.
\end{remark}

%==========================================================
\section{Устойчивость конструкций}\label{sec:stability}

Доказательства в этом разделе показывают, что требование (iii) из раздела~4 выполнено, и дают
явные константы. Это то, чего нет у теории концептуальных пространств: не «выпуклость
предположительно устойчива», а оценка.

\subsection{Липшицевость attention}

\begin{lemma}\label{lem:softmax}
Отображение $\softmax\colon\R^n\to\Delta^{n-1}$ липшицево с константой $1$ по
$\norm{\cdot}_\infty$, то есть $\norm{\softmax(u)-\softmax(u')}_\infty\le\norm{u-u'}_\infty$.
\end{lemma}
\begin{proof}
Компонента $\softmax_k(u)=e^{u_k}/Z(u)$, где $Z=\sum_j e^{u_j}$. Производная
$\partial\softmax_k/\partial u_j=\softmax_k(\delta_{kj}-\softmax_j)$; сумма модулей по $j$
равна $2\softmax_k(1-\softmax_k)\le 1/2$, а сумма квадратов не превосходит $\softmax_k(1-\softmax_k)$.
Отсюда градиент по $\norm{\cdot}_1$ не превосходит $1$, и утверждение следует из теоремы о
среднем значении.
\end{proof}

\begin{theorem}[Устойчивость внимания]\label{thm:attn}
Пусть $A^{(\ell)}$ и $\tilde A^{(\ell)}$ вычислены по векторам с $\norm{q_i}\le B$,
$\norm{k_j}\le B$. Тогда
\[
\norm{A^{(\ell)}-\tilde A^{(\ell)}}_\infty\ \le\ \frac{2B}{\sqrt d}\,\varepsilon,
\qquad \varepsilon=\max_i\norm{x_i^{(\ell)}-\tilde x_i^{(\ell)}} .
\]
\end{theorem}
\begin{proof}
Скалярное произведение билинейно и $1$-липшицево по каждому аргументу в sup-норме:
$\abs{q_i^\top k_j-\tilde q_i^\top\tilde k_j}\le \norm{q_i}\norm{k_j-\tilde k_j}+\norm{k_j}\norm{q_i-\tilde q_i}\le 2B\varepsilon$.
Деление на $\sqrt d$ и применение леммы~\ref{lem:softmax} покоординатно дают оценку.
\end{proof}

\subsection{Липшицевость спектра матрицы Грама}

\begin{theorem}[Неравенство Вейля]\label{thm:weyl}
Для симметричных $n\times n$ матриц $M,N$ собственные числа, упорядоченные по убыванию,
удовлетворяют $\max_k\abs{\lambda_k(M)-\lambda_k(N)}\le\norm{M-N}_F$.
\end{theorem}

\begin{theorem}[Устойчивость спектра]\label{thm:spec}
Пусть $\Sigma=V^\top V$, $\tilde\Sigma=\tilde V^\top\tilde V$. Тогда
\[
\max_k\abs{\lambda_k(\Sigma)-\lambda_k(\tilde\Sigma)}
\ \le\ \norm{\Sigma-\tilde\Sigma}_F
\ \le\ 2\sqrt n\,B\,\norm{V-\tilde V}_F .
\]
\end{theorem}
\begin{proof}
Первое неравенство — теорема~\ref{thm:weyl}. Второе:
$\norm{V^\top V-\tilde V^\top\tilde V}_F=\norm{V^\top(V-\tilde V)+(V-\tilde V)^\top\tilde V}_F\le \bigl(\norm{V}_F+\norm{\tilde V}_F\bigr)\norm{V-\tilde V}_F\le 2\sqrt n B\norm{V-\tilde V}_F$,
поскольку $\norm{V}_F\le\sqrt n\max_i\norm{v_i}\le\sqrt n B$.
\end{proof}

\begin{corollary}[Устойчивость эффективного ранга]
Отображение $\Sigma\mapsto R(\Sigma)$ непрерывно на множестве положительно определённых
матриц с нормированным следом; в частности, при $\lambda_k>\delta>0$ оно липшицево с константой,
зависящей от $\delta$ и $n$.
\end{corollary}

\subsection{Произведение метрик}

Из теорем~\ref{thm:attn} и~\ref{thm:spec} следует, что обе компоненты конструкции
$\Phi_{\mathrm{full}}$ устойчивы: матрица внимания — с константой $2B/\sqrt d$ по
$\norm{\cdot}_\infty$, спектр Грама — с константой $2\sqrt n\,B$ по норме
$\norm{\cdot}_F$. Для произведения метрик их можно объединить:
\[
L=\max\!\left(\frac{2B}{\sqrt d},\ 2\sqrt n\,B\right).
\]
Заметим, что из устойчивости $A$ устойчивость её спектра по неравенству Вейля следует с
константой $\abs{\lambda_k(A)-\lambda_k(\tilde A)}\le\norm{A-\tilde A}_F\le\sqrt n\norm{A-\tilde A}_\infty$;
численная константа при этом умножается на $\sqrt n$, и в приведённой оценке $L$ это учтено
через второй член. Это и есть количественная форма требования (iii). Обсуждение границ
применимости — в разделе~\ref{sec:limits}.

%==========================================================
\section{Эксперименты}\label{sec:experiments}

Все эксперименты делятся на шаг~0 (валидация конвейера) и четыре содержательных шага в порядке
выполнения. Шаг~0 предшествует остальным и проверяет, что измерительный прибор вообще работает и
даёт воспроизводимый сигнал (H1, раздел~\ref{sec:logic}); без него отрицательные результаты
неинтерпретируемы. Шаг~1 ставится \emph{первым} содержательным: он не требует эталона и
проверяет главное — устойчивость инварианта. Шаг~2 — вспомогательная проверка осмысленности
постановки. Шаг~3 — калибровка на размеченном корпусе; шаг~4 — топология. Шаги~3 и~4 имеют смысл
только после успеха шага~1.

\subsection{Шаг 0: валидация конвейера (H1)}\label{sec:h1}

Проводится до любых содержательных сравнений и ставит целью установить H1
(раздел~\ref{sec:logic}) — что $\Phi$ даёт воспроизводимый и неартефактный сигнал.

\begin{enumerate}[label=(\arabic*)]
\item \textbf{Тесты инвариантности (unit).} На одном тексте: (а)~детерминизм (двойной прогон
даёт бит-в-бит $\Phi$); (б)~инвариантность $\Phi_{\mathrm{spec}}$ к повороту $V\mapsto VR$ для
ортогональной $R$ (сингулярные числа сохраняются); (в)~инвариантность к перестановке позиций
$V\mapsto PV$ и $A\mapsto PAP^\top$. Это проверка реализации, а не гипотезы.
\item \textbf{Положительный контроль.} (i)~Синтетические $V$ с зашитым спектром
восстанавливаются $\Phi$; (ii)~известное реальное различение (число/согласование/часть речи)
отделяется. Не пройдено — конвейер невалиден, дальнейшие отрицательные результаты
неинтерпретируемы.
\item \textbf{Надёжность.} Расщепление выборки надвое (Спирмен—Браун, bootstrap-CI) и
репликация по независимым моделям/архитектурам. Сравнение — с достижимым потолком, а не с нулём.
\item \textbf{Конфаунды.} Partial Mantel против длины, частотности, синтаксических признаков;
проверка на выборке, выровненной по длине; базлайны (мешок слов, только длина).
\item \textbf{Выучено ли.} Untrained-контроль (случайная инициализация) и перемешивание порядка;
для $A$-компоненты сигнал обязан деградировать при перемешивании.
\item \textbf{Значимость и размер эффекта.} Перестановочные тесты; cross-validated AUC
«одно/разные»; выбор слоя/головы/метрики — на отложенной части (nested CV), иначе двойное
использование данных.
\end{enumerate}

Только после прохождения шага~0 имеет смысл шаг~1; провал шага~1 при пройденном шаге~0 —
содержательный отрицательный результат о гипотезе~\ref{hyp:relational}, а не о приборе.

\subsection{Шаг 1: устойчивость инварианта по слоям (без эталона)}

Это ключевой эксперимент, потому что он не требует знания, какие тексты задают одно понятие.
Достаточно ручного набора «заведомо одно» и «заведомо разное», а проверяется внутреннее
свойство: инвариант не должен «гулять» от слоя к слою.

\begin{criterion}[Стабильность инварианта по глубине]
Для текста $T$ и слоя $\ell$ вычисляется $\Phi^{(\ell)}(T)$ (по конструкции
$\Phi_{\mathrm{spec}}$ или $\Phi_{\mathrm{full}}$). Рассматривается последовательность
$\Phi^{(0)}(T),\Phi^{(1)}(T),\dots,\Phi^{(L)}(T)$. Критерий: для текстов одного понятия эта
последовательность «одинаково ведёт себя» (малые шаги между слоями, близкие значения на
поздних слоях у разных текстов), для текстов разных понятий — расходится.
\end{criterion}

\begin{criterion}[Разделимость по глубине]
Для каждой пары слоёв $(\ell,\ell')$ вычисляется отношение
\[
\frac{\delta_{\mathrm{out}}^{(\ell,\ell')}}{\delta_{\mathrm{in}}^{(\ell,\ell')}},
\qquad
\delta_{\mathrm{in}}^{(\ell,\ell')}=\max_{i,j}\norm{\Phi^{(\ell)}(T_i)-\Phi^{(\ell')}(T_j)}\ \ (T_i,T_j\ \text{одно понятие}),
\]
и аналогично $\delta_{\mathrm{out}}$ по разным понятиям. Если это отношение устойчиво
больше единицы на поздних слоях, инвариант отражает понятие; если оно ведёт себя как на
случайных текстах — нет.
\end{criterion}

\begin{remark}[Почему это первый эксперимент]
Он не зависит от LLM, не требует разметки корпуса и проверяет то, что гипотеза
\ref{hyp:training} предсказывает строго: если модель выучила связи понятия, они не должны
меняться от слоя к слою. Если сигнала нет здесь, ни шаг~3, ни LLM-калибровка его не
создадут. Отрицательный исход здесь — содержательный результат, указывающий на
опровержение гипотезы~\ref{hyp:training}.
\end{remark}

\subsection{Шаг 2: межслоевая согласованность (вспомогательный)}\label{sec:level0}

\begin{criterion}[Деформация относительно границы]
Для каждого слоя $\ell$ вычисляются
\[
D_A^{(\ell)}=\norm{A^{(\ell+1)}-A^{(\ell)}}_\infty,\qquad
D_\lambda^{(\ell)}=\max_k\abs{\lambda_k^{(\ell+1)}-\lambda_k^{(\ell)}},
\]
и отношения
\[
\rho_A^{(\ell)}=\frac{D_A^{(\ell)}}{(2B/\sqrt d)\,\varepsilon^{(\ell)}},\qquad
\rho_\lambda^{(\ell)}=\frac{D_\lambda^{(\ell)}}{2\sqrt n\,B\,\varepsilon^{(\ell)}} ,
\]
где $\varepsilon^{(\ell)}=\max_i\norm{x_i^{(\ell+1)}-x_i^{(\ell)}}$.
\end{criterion}

Интерпретация: $\rho\ll 1$ означает, что преобразование слоя деформирует структуру
\emph{меньше}, чем допускает худший случай, то есть структура сохраняется сверх гарантированного;
$\rho\approx 1$ означает, что деформация упирается в границу, и межслоевой инвариантности нет.

\begin{remark}[Что этот эксперимент \emph{не} измеряет]
Шаг~2 измеряет \emph{межслоевую} согласованность и ничего не говорит о существовании
островов, которые являются свойством \emph{внутрислойного} распределения представлений
текстов. Положительный исход здесь — условие необходимости, а не достаточности, и
формулировать его как подтверждение гипотезы об островах было бы некорректно.
\end{remark}

\subsection{Шаг 3: прямой тест на размеченном корпусе}\label{sec:level2}

Проводится только при успехе шага~1. Здесь наконец нужен эталон, и его роль разбирается в
разделе~\ref{sec:nogroundtruth}: ручной набор «одно / разные понятия», при расширении —
LLM-разметка корпуса без права быть судьёй.

\paragraph{Изоморфизм структур (RSA).} Помимо бинарного разделения «одно/разные»,
содержательнее проверить выравнивание \emph{структур второго порядка}
(раздел~\ref{sec:logic}): построить RDM инвариантов и RDM независимых человеческих оценок
сходства на одном наборе понятий и сравнить ранговой корреляцией. Кандидаты
$\Phi_{\mathrm{pool}}$, $\Phi_{\mathrm{spec}}$, $\Phi_{\mathrm{full}}$ сравниваются на этом
едином критерии: конструкция, чья RDM лучше согласуется с человеческой, и есть эмпирически
выбранная (это снимает открытый вопрос выбора канонизации из раздела~\ref{sec:order}).
Предсказание: $\Phi_{\mathrm{pool}}$ уступает, $\Phi_{\mathrm{spec}}/\Phi_{\mathrm{full}}$ —
нет. Предварительный прогон (раздел~\ref{sec:prelim}) дал обратное: выигрывает направление,
а $\Phi_{\mathrm{spec}}$ смешана с длиной.

Выбираются понятия $P_1,\dots,P_m$ и для каждого — $k$ текстов, задающих это понятие
(например, определение, синонимичное описание, пример употребления). Вычисляются
\begin{equation}\label{eq:crit}
\delta_{\mathrm{in}}=\max_{i,j}\norm{\Phi(T_i)-\Phi(T_j)},\qquad
\delta_{\mathrm{out}}=\min_{T\in\mathcal T_{P_i},\,S\in\mathcal T_{P_j},\,i\ne j}\norm{\Phi(T)-\Phi(S)}.
\end{equation}

\begin{criterion}[Остров]
Понятие обладает компактным представлением, если $\delta_{\mathrm{in}}\ll\delta_{\mathrm{out}}$.
Отношение $\delta_{\mathrm{out}}/\delta_{\mathrm{in}}$ измеряет \emph{чёткость} понятия.
\end{criterion}

\begin{criterion}[Непрерывность]
Для пары текстов одного понятия, отличающихся малым изменением, вычисляется
\[
\eta=\max_{i}\norm{\Phi(T_i)-\Phi(T_i')}_{\text{малое изменение}}.
\]
Гипотеза~\ref{hyp:continuity} предсказывает, что $\eta$ мала в масштабе $\delta_{\mathrm{in}}$.
\end{criterion}

\paragraph{План измерений.}
Три группы текстов на одно понятие по возрастанию длины: одно слово; короткое определение;
развёрнутое описание (абзац). Для каждой группы вычисляются $\delta_{\mathrm{in}}$,
$\delta_{\mathrm{out}}$, $\eta$ и $R(\Sigma)$. Ожидаемая динамика: $R$ падает с ростом текста
(если верен тезис «длиннее — конкретнее»); если $R$ не падает, тезис опровергнут, и это само по
себе содержательный результат.

\paragraph{Ожидаемые исходы.}
\begin{enumerate}[label=(\alph*)]
\item $\delta_{\mathrm{in}}\ll\delta_{\mathrm{out}}$ устойчиво и $\eta$ мала: гипотезы
подтверждаются, понятие действительно стягивается в отделённый остров.
\item $\delta_{\mathrm{in}}$ сравнима с $\delta_{\mathrm{out}}$: представления отдельных текстов
перемешаны, понятие не стягивается; необходимо пересмотреть $\Phi$.
\item $\eta$ велика при малых изменениях текста: непрерывность нарушается, и тогда сам выбор
пространства представлений неудачен.
\end{enumerate}

\paragraph{Согласование с шагом~1.} Результат шага~3 имеет вес только если согласуется с
шагом~1: если инвариант стабилен по глубине, но не разделяет понятия на размеченном корпусе,
это указывает либо на дефект корпуса, либо на то, что понятие в модели не локализовано как
область; различить эти два случая — отдельная задача.

\subsection{Шаг 4: топология островов}\label{sec:topology}

Проводится только при успехе шагов~1 и~3. На облаке $\{\Phi(T)\}$ для фиксированного слоя
строится фильтрация Вьеториса—Рипса и вычисляется персистентная гомология. Значимыми считаются
классы, живущие при радиусах, сравнимых с $\delta_{\mathrm{out}}$, но существенно превосходящих
$\delta_{\mathrm{in}}$. Долгоживущие циклы $H_1$ означают наличие «провалов» между островами —
участков пространства представлений, не занятых ни одним понятием.

\begin{remark}[Опасность «структуры ради структуры»]
Вычисление гомологий полезно только при наличии независимой гипотезы о том, что означает
топологический признак. Если такой гипотезы нет, результат сводится к переописанию уже
известного на другом языке — именно та претензия, которая справедливо адресована большой части
работ по топологическому анализу представлений. Поэтому шаг~4 стоит после шагов~1 и~3, а
не вместо них.
\end{remark}

%==========================================================
\section{Предварительные эмпирические результаты}\label{sec:prelim}

Ниже — что дал предварительный прогон шагов~0--3 на малой модели
(\texttt{distilbert-base-uncased} для двунаправленного случая и \texttt{distilgpt2}
для причинного; CPU) на небольшом ручном корпусе (парафразы конкретных
существительных). Числа приводятся как ориентир, а не как установленный факт: корпус
мал, модели и язык единичны. Но направление результата устойчиво и существенно для
программы.

\subsection{Шаг 0: спектр Грама не проходит проверку H1}

Отпечаток $\Phi_{\mathrm{spec}}(T)$ (нормированный спектр $\tilde\Sigma$) на уровне
целого предложения оказался по существу мерой \emph{длины}: в слоях~1--5 геометрия
инвариантов объяснялась числом токенов на $72$--$94\%$ (коэффициент детерминации), а
коэффициент надёжности (расщепление пополам, поправка Спирмена—Брауна) был около нуля
или отрицателен. Разделение «одно/разные понятия» давало AUC $\approx0.53$--$0.60$
(случайное угадывание — $0.5$), а после вычета пересечения слов и длины падало до
$\approx0.49$. \emph{Причина структурна}: при $n<d$ число ненулевых сингулярных чисел
$V$ равно $n$ (ранг $\le n$), поэтому сам носитель спектра кодирует число токенов;
нормировка на след этого не убирает. Положительный контроль (синтетические $V$ с
заложенным спектром) проходил (AUC $=1.0$, надёжность $\approx0.86$), то есть прибор
исправен, но измеряет не то.

\subsection{Направление против соотношений}

В тех же условиях семантический сигнал нёс усреднённый вектор (направление), а не
спектр. На парафразах с почти непересекающимся словарём и выровненной длиной
(содержательное пересечение слов $\approx0.006$; лексический базлайн — AUC $0.52$, то
есть слова не разделяют):
\begin{center}
\begin{tabular}{@{}lccc@{}}
\toprule
сигнал & AUC & AUC после контроля & AUC при нулевом пересечении\\
\midrule
$\Phi_{\mathrm{spec}}$ (спектр Грама) & $0.53$ & $0.60$ & $0.53$\\
усреднённое направление + косинус & $0.996$ & $0.99$ & $0.995$\\
совместный вход (нативное сравнение) & $1.00$ & $0.98$ & $1.00$\\
\bottomrule
\end{tabular}
\end{center}
То же на нативном сравнении: косинус отдельно закодированных усреднённых векторов
воспроизводил сравнение двух текстов «в одном входе» с корреляцией $\approx0.88$, а
выученная на направлениях метрика — с $\approx0.93$; $\Phi_{\mathrm{spec}}$ давала
$\approx0.15$--$0.18$. Иначе говоря, экспорт «работает», но носителем оказывается
направление, которое спектр Грама \emph{по построению отбрасывает}.

\subsection{Родственные понятия: грубое против тонкого}

На корпусе родственных пар (пчела/оса, река/ручей, снег/град, вулкан/гейзер,
гитара/скрипка, акула/дельфин) с непересекающимися словами:
\begin{center}
\begin{tabular}{@{}lccc@{}}
\toprule
контраст & $\Phi_{\mathrm{spec}}$ & направление & нативное\\
\midrule
same vs related (тонко) & $0.47$ & $0.64$--$0.67$ & $0.65$\\
related vs unrelated (грубо) & $0.57$ & $0.96$ & $0.96$\\
порядок same$<$related$<$unrelated, $\rho$ & $-0.05$ & $-0.69$ & $-0.71$\\
\bottomrule
\end{tabular}
\end{center}
Направление даёт правильную \emph{грубую} структуру (родственное ближе несвязанного) и
почти идеально разделяет «родственное vs несвязанное», но \emph{тонкое} различение
понятия и его соседа (пчела/оса) ему почти не даётся: лучший результат по слоям — AUC
$\approx0.67$ (слой~5 из~6), а последний слой для тонкого даже хуже. Спектр Грама не
даёт ни того, ни другого.

\subsection{Траектория и аттрактор: понятие как транзиент}

Отдельная серия опытов била не в мгновенное состояние, а в \emph{динамику} генерации:
давали описание понятия, затем один и тот же вопрос и смотрели короткую траекторию
(3--6 токенов), а также распределение токена сразу после вопроса.

Ключевое наблюдение: непосредственно следующее распределение (<<запуск>> траектории)
различает понятия, а сама траектория быстро уходит в общий шаблон. На корпусе родственных
пар (вопрос вида <<It is a>>, модель \texttt{Qwen2.5-1.5B-Instruct}):
\begin{center}
\begin{tabular}{@{}lcc@{}}
\toprule
сигнал & same vs related & related vs unrelated\\
\midrule
распределение первого токена (JS) & $0.67$ & $0.92$\\
эмбеддинги траекторий по 6 токенов & $0.48$ & $0.60$\\
\bottomrule
\end{tabular}
\end{center}
То есть запуск даёт лучшее тонкое разделение за всю серию ($0.67$ против $0.65$ у входного
направления и $\approx0.5$ у косинуса слов), но как только генерация делает несколько
шагов, все понятия сваливаются в один и тот же \emph{стилевой} аттрактор вида <<It is a
fact that\ldots>>, <<It is a type of\ldots>>, <<It is a very common\ldots>>, и эмбеддинги
траекторий становятся неразличимы.

Отсюда уточнение: \emph{аттрактор динамики существует, но он общий, а не понятийный;
понятие живёт во входе в аттрактор (в транзиенте), а не в длинном пределе}. Это
согласуется с механизмом: генерация есть повторное применение фиксированного отображения
<<префикс $\to$ распределение следующего токена>> (с KV-кэшем, так что каждый шаг обусловлен
всем префиксом). Высоковероятные общие продолжения (<<fact that>>, <<type of>>) образуют
глубокий бассейн, независимый от темы; попав туда за один-два токена, модель теряет
понятийную специфику.

Отдельно зафиксируем разрыв, обнаруженный прямыми пробами: модель
\texttt{Qwen2.5-1.5B} \emph{знает} разницу пчела/оса (в свободном ответе перечисляет
различия), но косинус и внимание между векторами этих двух слов её \emph{не} отражают
(<<пчела/оса>> и <<пчела/вулкан>> неразличимы). Значит \textbf{отношение не есть близость
объектов отношения}: знание хранится процедурно, а не как расстояние между двумя точками.
(Модель 0.5B при этом вообще проваливает контроль знания и для выводов о представлениях
непригодна.)

\subsection{Что из этого следует для программы}

Предварительный результат \emph{противоречит} исходной ставке на соотношения: на этом
материале основной сигнал сидит в направлении, а не в инвариантной к повороту форме, и
спектр Грама отбрасывает именно то, что информативно. Это не опровержение гипотезы о
понятиях, а опровержение одного конкретного кандидата в отпечаток
($\Phi_{\mathrm{spec}}$) и его ранжирования относительно $\Phi_{\mathrm{pool}}$. Тонкое
различение (пчела/оса) не даётся ни усреднённым направлением, ни спектром; вероятное
место его поиска — токенный уровень и паттерны внимания (усреднение размывает признаки).

Ограничения результата существенны: корпус мал и ручен; понятия — конкретные
существительные; проверены одна-две модели, английский и сжатый набор слоёв;
кросс-длинный режим (контрпример с пчелой, раздел~1.4) отдельно не проверялся, и там
направления ещё могут сломаться; человеческий эталон (H2) не подключался. Поэтому вывод
— не «направления побеждают вообще», а «на данном материале $\Phi_{\mathrm{spec}}$ не
работает, а направление даёт грубый, но воспроизводимый сигнал; тонкое различение
остаётся открытым».

Совокупность результатов меняет постановку цели. Искать геометрический инвариант понятия
(точку, отпечаток, остров в пространстве векторов) на проверенном материале не удаётся:
спектр Грама смешан с длиной, усреднённое направление ловит лишь грубую категорию, а
тонкое различие (пчела/оса) не лежит ни в косинусе, ни во внимании даже у модели, которая
это различие знает. Наиболее правдоподобная картина: \emph{понятие хранится процедурно} —
как транзиент генерации и как умение, разворачивающееся в последовательности, — а не как
сравнимое место в пространстве. Поэтому осмысленная цель сужается до двух: (а)
\emph{поведенческая} — какие различения модель делает надёжно и по-человечески; (б)
\emph{механистическая} — где и как данное различение вычисляется (транзиент/запуск,
пробинг, причинные вмешательства). Человеческий эталон (H2) при этом остаётся
незакрытым: без него любое утверждение о совпадении с человеческими понятиями
необоснованно.

%==========================================================
\section{Границы применимости и что результаты \emph{не} означают}\label{sec:limits}

\paragraph{Слабость линейных оценок.} Теоремы~\ref{thm:attn} и~\ref{thm:spec} дают оценку
худшего случая, растущую с числом позиций $n$. Для длинных текстов линейная граница грубеет.
Кроме того, $\norm{V-\tilde V}_F$ измеряет \emph{поэлементное} возмущение, тогда как реальная
контекстуализация меняет \emph{структуру} смешивания, а не только амплитуды. Линейные оценки —
верхняя грань, а не описание наблюдаемого поведения.

\paragraph{Модельная зависимость.} Все результаты относятся к конкретной архитектуре и весам.
«Структура устойчива в BERT» не влечёт «структура устойчива в человеческом мышлении» и вообще
не влечёт межмодельной инвариантности. Трансформер — техническая реализация, и его особенности
не следует переносить на познание без отдельного аргумента.

\paragraph{Что подтверждает шаг 1.} Что инвариант (спектр Грама и/или спектр внимания)
стабилен по слоям у текстов одного понятия и расходится у разных. Это внутреннее свойство
модели, и оно не требует эталона — но и не является доказательством, что инвариант
\emph{выражает} понятие: он может отражать устойчивое свойство поверхностного уровня
(длину, синтаксис, частотность), а не содержательное единство.

\paragraph{Что подтверждает шаг 3.} Что тексты одного понятия стягиваются в компактную
область на размеченном корпусе. Это эмпирический факт о данной модели и данном наборе текстов,
а не теорема о природе понятий. Из-за отсутствия независимого эталона сила вывода ограничена
качеством разметки.

\paragraph{Что подтверждает шаг 4.} Что между областями есть топологический дефект
определённого рода. Объяснения того, почему одни понятия образуют дыры, а другие нет, топология
не даёт — это вопрос к теории содержания понятий, а не к её геометрии.

%==========================================================
\section{Связь с исходной рамкой}

Программа не требует ответа на вопрос о критерии тождества, но и не подменяет его. Если
гипотезы раздела~5 подтверждаются, получается: представление текста стягивается в область и
изменяется непрерывно, — то есть представление ведёт себя как объект с устойчивой геометрией.
Это ещё не тождество: близость и компактность не дают правила «то же самое или нет». То есть
геометрия здесь обеспечивает \emph{условие осмысленности}, а не сам критерий. Именно поэтому
программа может быть выполнена независимо от того, принимается ли различение
прото-символа и символа: она даёт проверяемый слой, поверх которого это различение
формулируется, а не заменяет его собой.

Отдельно подчеркнём, что программа \emph{не} утверждает: что представление понятия — это
направление вектора, что понятие геометрически вложено в сферу слоя, или что свойства, найденные
в трансформере, суть свойства человеческого мышления. Первое опровергнуто контрпримером
раздела~1.4 и выведено из конструкций; второе не следует из обучения; третье — отдельная
гипотеза, для которой в документе нет оснований.

%==========================================================
\section*{Приложение: сводка величин}
\addcontentsline{toc}{section}{Приложение: сводка величин}

\begin{table}[h]
\centering
\begin{tabular}{@{}lll@{}}
\toprule
Величина & Определение & Что измеряет\\
\midrule
$A$ & $\softmax_j(q_i^\top k_j/\sqrt d)$ & граф полезных для предсказания связей\\
$\Sigma$ & $V^\top V$ & матрица Грама (все соотношения векторов)\\
$\tilde\Sigma$ & $\Sigma/\tr\Sigma$ & нормировка, снимающая зависимость от длины\\
$R(\Sigma)$ & $\exp(-\sum p_k\ln p_k)$ & эффективный ранг (конкретность)\\
$\Phi_{\mathrm{pool}}$ & пул векторов & базовая линия (контроль)\\
$\Phi_{\mathrm{spec}}$ & $\spec(\tilde\Sigma)$ & инвариант соотношений, без направления\\
$\Phi_{\mathrm{full}}$ & $(\spec(A),\spec(\tilde\Sigma))$ & соотношения плюс структура связей\\
$D_A$ & $\norm{A^{(\ell+1)}-A^{(\ell)}}_\infty$ & межслоевая деформация внимания\\
$D_\lambda$ & $\max_k\abs{\lambda_k^{(\ell+1)}-\lambda_k^{(\ell)}}$ & межслоевая деформация спектра\\
$\rho_A,\rho_\lambda$ & отношения к границам теорем~\ref{thm:attn},~\ref{thm:spec} & запас устойчивости\\
$\delta_{\mathrm{in}},\delta_{\mathrm{out}}$ & формула~\eqref{eq:crit} & отделимость понятий\\
$\eta$ & $\max_i\norm{\Phi(T_i)-\Phi(T_i')}$ & непрерывность при малом изменении\\
$\sigma_i(V)$ & сингулярные числа $V$ & вход в $\Phi_{\mathrm{spec}}$\\
$p$ & $\sigma_i^2/\sum_j\sigma_j^2$ & спектральное распределение (вероятностный вектор)\\
$d_H$ & $\sqrt{1-\sum_k\sqrt{p_kq_k}}$ & расстояние Хеллингера между $p$\\
$\mathrm{rel}$ & корреляция двух половин / моделей & надёжность измерения (H1)\\
$\rho_\text{partial}$ & частная корреляция (Мантел) & сигнал сверх конфаундов\\
$\mathrm{AUC}$ & cross-validated & разделимость «одно/разные»\\
\bottomrule
\end{tabular}
\caption{Сводка используемых величин.}
\end{table}

\end{document}
